% concatenated from 110.tex
\documentclass[11pt]{amsart}

\usepackage{amsmath,amssymb,amsthm,mathtools}
\usepackage{booktabs}
\usepackage{enumitem}
\usepackage[hidelinks]{hyperref}
\usepackage{url}
\usepackage[margin=1.05in]{geometry}

\newtheorem{theorem}{Theorem}[section]
\newtheorem{proposition}[theorem]{Proposition}
\newtheorem{lemma}[theorem]{Lemma}
\newtheorem{corollary}[theorem]{Corollary}
\theoremstyle{remark}
\newtheorem{remark}[theorem]{Remark}

\newcommand{\F}{\mathbb F}
\newcommand{\Aut}{\operatorname{Aut}}
\newcommand{\Syl}{\operatorname{Syl}}
\newcommand{\PP}{\mathbb P}
\newcommand{\cH}{\mathcal H}

\title{Transitive automorphism groups of maximal curves}

\author{Saeed Tafazolian}
\address{Institute of Mathematics, Statistics and Scientific Computing, University of Campinas, Brazil}
\email{saeed@unicamp.br}

\subjclass[2020]{11G20, 14G15, 14H37}
\keywords{maximal curve, Hermitian curve, automorphism group, Klein quartic, transitive action}

\begin{document}

\begin{abstract}
Let $q=p^h$ be an odd prime power and let $X/\F_{q^2}$ be a maximal
curve of genus at least two. We classify the curves for which the full
geometric automorphism group is transitive on the set of
$\F_{q^2}$-rational points. We prove that, if $h>1$, then $X$ is the
Hermitian curve. If $h=1$, the only additional possibility occurs for
$q=5$: the unique $\F_{25}$-maximal genus-three curve, namely the
maximal $S_4$-model of the Klein quartic. The proof separates tame and
wild actions. In the tame case, genus bounds, signatures of quotient
maps, specialization to characteristic zero, low-genus automorphism
classifications, and a Hasse--Witt computation reduce the problem to
the Klein quartic. In the wild case, the rational points are identified
with the Sylow $p$-subgroups of the automorphism group. Noncyclic Sylow
subgroups are treated using a finite-group classification theorem
together with Henn's large-automorphism classification, while cyclic
Sylow subgroups are excluded by local ramification and the
Riemann--Hurwitz formula. Combined with the known characteristic-two
result, this yields the classification for all prime powers.
\end{abstract}

\maketitle

\section{Introduction}

Let $q=p^h$ be a power of an odd prime and let $X/\F_{q^2}$ be a
projective, nonsingular, geometrically irreducible curve of genus
$g\ge2$.  Recall that $X$ is maximal if
\[
        \#X(\F_{q^2})=q^2+1+2gq.
\]
We are interested in the case in which the full geometric automorphism
group of $X$ is transitive on $X(\F_{q^2})$.

The basic example is the Hermitian curve
\[
        \cH_{q+1}:\qquad y^{q+1}=x^q+x.
\]
It has genus $q(q-1)/2$ and $q^3+1$ rational points, and
$\Aut(\cH_{q+1})\cong PGU(3,q)$ acts doubly transitively on
$\cH_{q+1}(\F_{q^2})$; see \cite[Theorem~12.24(iv)]{HKT}.  In odd
characteristic there is one further example: the Klein quartic has an
$\F_{25}$-maximal model of genus $3$ with $56$ rational points, and its
full automorphism group of order $168$ is transitive on these points.

The transitivity problem was considered by Giulietti, Korchm\'aros and
Montanucci in \cite[Section~6.1]{GKMSurvey}.  Among the standard
families of maximal curves, the Hermitian curve is distinguished by the
fact that its full automorphism group is transitive on all rational
points.  In characteristic $2$, Giulietti and Korchm\'aros proved that
this property characterizes the Hermitian curve among maximal curves of
genus greater than one; see \cite[Theorem~9.1]{GK2010}.  The purpose of
this paper is to settle the odd-characteristic case.

Our main result is the following.

\begin{theorem}\label{thm:main}
Let $q=p^h$ be an odd prime power and let $X/\F_{q^2}$ be a maximal
curve of genus $g\ge2$.  Then $\Aut(X)$ acts transitively on
$X(\F_{q^2})$ if and only if one of the following holds:
\begin{enumerate}[label=\textup{(\roman*)}]
\item $X$ is $\F_{q^2}$-isomorphic to the Hermitian curve
\[
             \cH_{q+1}:\qquad y^{q+1}=x^q+x;
\]
\item $q=5$ and $g=3$. Equivalently, $X$ is the unique
$\F_{25}$-maximal curve of genus $3$, represented by the classical
``$S_4$-model'' of the Klein quartic introduced by Elkies
\cite[(1.11)]{ElkiesKlein}; see also \cite[\S3.1]{BTT} for its
$\F_{25}$-maximal realization.
\end{enumerate}
\end{theorem}

We divide the proof into the tame and wild cases.  In the tame case we
prove a slightly stronger statement, allowing an arbitrary subgroup of
order prime to $p$ which is transitive on $X(\F_{q^2})$.  A numerical
reduction first gives $q=p\le19$ for a non-Hermitian curve.  The upper
part of the genus spectrum then leaves seven possibilities.  They are
excluded by tame specialization, low-genus automorphism bounds,
triangle-group arguments and a Hasse--Witt computation, with the
single exception of the maximal Klein quartic over $\F_{25}$.

In the wild case, maximality implies that $X$ has $p$-rank zero.  The
rational points are then naturally in bijection with the Sylow
$p$-subgroups of the full automorphism group; in particular,
\[
   |\Syl_p(\Aut(X))|=q^2+1+2gq.
\]
For a noncyclic Sylow $p$-subgroup we use the group-theoretic theorem of
Guralnick--Malmskog--Pries together with Henn's classification.  The
former uses the classification of finite simple groups; this is the
only place in the proof where we invoke CFSG directly.  The cyclic case
is treated separately by combining the global Riemann--Hurwitz formula
with the local ramification filtration.

The hypothesis of Theorem~\ref{thm:main} is different from the usual
``many automorphisms'' assumptions.  For $q=p$, \cite{BMT} obtains
Hermitian coverings under a lower bound on the size of the automorphism
group, while \cite{GK2019} studies large automorphism groups in relation
to the $p$-rank.  Here no order threshold is imposed; the transitivity
assumption itself gives the orbit structure used in the proof.

Throughout the paper, $\Aut(X)$ denotes the geometric automorphism
group.  Every geometric automorphism of a maximal curve over
$\F_{q^2}$ is defined over $\F_{q^2}$; see \cite[Theorem~3.10]{GSY}.
By ``the Hermitian curve'' over $\F_{q^2}$ we mean the
$\F_{q^2}$-isomorphism class of $\cH_{q+1}$.  When a curve is only
geometrically isomorphic to $\cH_{r+1}$ for some power $r$ of $p$,
possibly with $r\ne q$, we shall say explicitly that it is
geometrically Hermitian.

\section{Preliminaries}

We collect the facts used in the proof.

\subsection{Genus bounds}

We shall repeatedly use the upper part of the genus spectrum of a
maximal curve.  The maximal-genus characterization in part \textup{(i)}
is due to R\"uck--Stichtenoth, while the second-genus bound in part
\textup{(ii)} is due to Fuhrmann--Garcia--Torres; we use the formulations
recalled in \cite[Lemma~2.12 and Corollary~2.13]{KT}.  The equality case
in part \textup{(iii)} is \cite[Theorem~3.1]{FGT}, and the third genus
bound in part \textup{(iv)} is \cite[Corollary~3.3]{KT}.

\begin{lemma}\label{lem:gaps}
Let $X/\F_{q^2}$ be a maximal curve of genus $g$.
Then:
\begin{enumerate}[label=\textup{(\roman*)}]
\item
$g\le q(q-1)/2$, with equality if and only if $X$ is Hermitian.

\item If $X$ is not Hermitian, then
\[
g\le \frac{(q-1)^2}{4}.
\]

\item If $q$ is odd and $g=(q-1)^2/4$, then
\[
X\simeq_{\F_{q^2}}\{\,y^{(q+1)/2}=x^q+x\,\}.
\]

\item If
\[
g<\left\lfloor\frac{(q-1)^2}{4}\right\rfloor,
\]
then
\[
g\le\left\lfloor\frac{q^2-q+4}{6}\right\rfloor.
\]
\end{enumerate}
\end{lemma}

A maximal curve is supersingular.  We shall also use the following
standard consequence of the Kleiman--Serre covering result; see
\cite[Proposition~6]{Lachaud}.

\begin{lemma}\label{lem:quotientmax}
Let $X/\F_{q^2}$ be maximal, and let $H\leq\Aut(X)$ be defined over
$\F_{q^2}$.  Then the quotient $X/H$ is maximal over $\F_{q^2}$.
In particular, every positive-genus quotient of $X$ is supersingular.
\end{lemma}

\subsection{Group actions and ramification notation}

Let a finite group $G$ act on a curve $X$.  For $P\in X$, we write
\[
        G_P=\{\sigma\in G:\sigma(P)=P\}
\]
for the stabilizer of $P$, and
\[
        G\cdot P=\{\sigma(P):\sigma\in G\}
\]
for its orbit.  The orbit--stabilizer formula gives
\[
        |G|=|G_P|\,|G\cdot P|.
\]
An orbit is called short if its length is smaller than $|G|$, or
equivalently if its point stabilizers are nontrivial.

The action determines the quotient Galois cover
\[
        X\longrightarrow X/G.
\]
For $P\in X$, the ramification index at $P$ is $|G_P|$.  We write
\[
        G_P=G_P^{(0)}\supseteq G_P^{(1)}
             \supseteq G_P^{(2)}\supseteq\cdots
\]
for the lower ramification groups.  If the action is tame, then
$G_P^{(1)}=1$ and $G_P$ is cyclic.  In Section~\ref{sec:tame} we shall
call a subgroup $G\le\Aut(X)$ \emph{tame} when $p\nmid|G|$.  This is
the global hypothesis used there; in particular every inertia group of
the action is tame.

We write $\Syl_p(G)$ for the set of Sylow $p$-subgroups of $G$,
$N_G(S)$ for the normalizer of a subgroup $S\le G$, and
$O_{p'}(G)$ for the largest normal subgroup of $G$ of order prime to
$p$.  A Sylow $p$-subgroup is called TI if distinct conjugates have
trivial intersection.

We shall use the following structural result in the wild case. It is
quoted here in the geometric form needed below, in order to make clear
both its hypotheses and its $2$-transitivity conclusion.

\begin{theorem}[Guralnick--Malmskog--Pries in the form used here]\label{thm:GMP-TI}
Let $A$ act faithfully on a curve $Y$ of genus at least two over an
algebraically closed field of characteristic $p$, and let
$Q\in\Syl_p(A)$. Suppose that $Q$ is not cyclic (and, when $p=2$, not
generalized quaternion), that $Q$ fixes exactly one point of $Y$ and
acts freely on its complement, and that
\[
        N_A(Q)=QC
\]
with $C$ cyclic. If $N_A(Q)\ne A$ and $M=Z(N_A(Q))$, then $M$ is a
normal $p'$-subgroup of $A$, the quotient $A/M$ is almost simple, and
its socle is isomorphic to one of
\[
 PSL(2,p^a)\quad(a\ge2),\qquad
 PSU(3,p^a)\quad(p^a>2),
\]
\[
 Sz(2^{2a+1})\quad(p=2,\ a>1),\qquad
 {}^2G_2(3^{2a+1})'\quad(p=3).
\]
Moreover, $A$ acts doubly transitively on $\Syl_p(A)$.
\end{theorem}

\begin{proof}[Justification of the quoted form]
The normality of $M$ and the almost-simple conclusion are precisely the
geometric conclusion of \cite[Theorem~3.19]{GMP}; we record why the
$2$-transitivity assertion from \cite[Theorem~3.16]{GMP}, on which that
theorem is based, applies in the present setting.

First, $Q$ is a TI subgroup. Let $P$ be the unique fixed point of $Q$
and take $a\notin N_A(Q)$. If $Q$ and $Q^a$ fixed the same point, then
both would be Sylow $p$-subgroups of the corresponding point
stabilizer. The first ramification group is the unique Sylow
$p$-subgroup of a point stabilizer, so this would force $Q=Q^a$, a
contradiction. Thus the fixed points of $Q$ and $Q^a$ are distinct. A
nontrivial element of $Q\cap Q^a$ would fix both of them, contrary to
the assumed free action away from the unique fixed point of $Q$. Hence
$Q$ is TI.

There is no mismatch between the hypotheses used here and those in
\cite{GMP}. In Notation~3.18 of that paper the condition that $Q$ be
noncyclic (and, for $p=2$, not generalized quaternion) is explicitly
noted to be equivalent to the existence of an elementary abelian
subgroup of order $p^2$. Moreover, Theorem~3.19 assumes
$A\ne I=N_A(Q)$, which is exactly our hypothesis $N_A(Q)\ne A$.
The group-theoretic Theorem~3.16 is stated under the weaker condition
$Q\ne A$, and this is automatic here.

For completeness, we make explicit the initial reduction in the proof
of \cite[Theorem~3.16]{GMP}. Since $N_A(Q)\ne A$, choose
$a\notin N_A(Q)$. The normal $p$-subgroup $O_p(A)$ is contained in
every Sylow $p$-subgroup of $A$, and therefore
\[
        O_p(A)\le Q\cap Q^a=1.
\]
Thus $O_p(A)=1$. The argument of \cite[Theorem~3.16]{GMP} then passes
to $A/O_{p'}(A)$, applies the finite-simple-group classification quoted
there, and obtains the displayed socle list and the doubly transitive
action on the Sylow $p$-subgroups. In fact, that theorem proves the
slightly stronger identity $M=O_{p'}(A)$. The classification of finite
simple groups enters only through this step; see
\cite[Remark~3.17]{GMP}.
\end{proof}

\subsection{Tame specialization and triangle actions}\label{subsec:tame-specialization}

Suppose that $p\nmid|G|$ and that the quotient $X/G$ is rational with
exactly three branch points.  If the ramification indices are
$(a,b,c)$, with
\[
       2\le a\le b\le c,\qquad
       \frac1a+\frac1b+\frac1c<1,
\]
then Grothendieck's specialization theorem for the prime-to-$p$
fundamental group of the punctured line gives a generating triple of $G$
of orders $a,b,c$ whose product is $1$; see
\cite[Expos\'e~XIII, Corollaire~2.12]{SGA1}.  Equivalently, $G$ is a
finite quotient of the hyperbolic triangle group $\Delta(a,b,c)$.  By the
Riemann existence theorem this same generating triple gives a complex
$G$-cover of $\PP^1$ with the same signature, and hence, by
Riemann--Hurwitz, with the same genus.  This is the only specialization
consequence used below.

The tame Riemann--Hurwitz formula reads
\begin{equation}\label{eq:triangleRH}
  2g-2
   =|G|\left(1-\frac1a-\frac1b-\frac1c\right).
\end{equation}

\begin{remark}\label{rem:threebranch}
We shall repeatedly use the following elementary consequence of the tame
Riemann--Hurwitz formula.  Let $G\le\Aut(X)$ satisfy $p\nmid|G|$. If
\[
        \frac{|G|}{g-1}>12,
\]
then $X/G$ is rational and the cover $X\to X/G$ has exactly three branch
points.  Indeed, if $X/G$ has positive genus, Riemann--Hurwitz gives
$|G|/(g-1)\le4$.  If $X/G\simeq\PP^1$ has at least four branch points,
then $|G|/(g-1)\le12$; for four branch points the smallest positive
orbifold defect is $1/6$, attained by $(2,2,2,3)$, and for at least five
branch points the defect is at least $1/2$.
\end{remark}

All characteristic-zero actions used below have genus at most $48$.
We use Breuer's classification of finite groups of holomorphic
(equivalently, orientation-preserving) automorphisms of compact Riemann
surfaces \cite[Chapter~5 and the summary table on p.~91]{Breuer}.
We shall also use Burnside's normal $p$-complement theorem and his
$p^a q^b$ theorem in the forms stated in
\cite[Theorems~5.13 and~7.8]{Isaacs}.

\subsection{\texorpdfstring{The curves $\cH_m$}{The curves Hm}}

For later use in the prime-field case, let $m$ be a divisor of $p+1$ and put
\[
            \cH_m:\qquad y^m=x^p+x .
\]
Then
\begin{equation}\label{eq:Hmgenus}
            g(\cH_m)=\frac{(m-1)(p-1)}2.
\end{equation}
For $2\le m<p+1$ one has
\begin{equation}\label{eq:Hmaut}
            |\Aut(\cH_m)|=m\,p(p^2-1);
\end{equation}
indeed $\Aut(\cH_m)$ contains a normal cyclic group of order $m$ with
quotient $PGL(2,p)$; see \cite[Lemma~2.3]{BMT}.


\section{The tame case}\label{sec:tame}

Throughout this section $G\le\Aut(X)$ is a subgroup of order prime to
$p$ which acts transitively on $X(\F_{q^2})$.  We first isolate the
numerical range in which such an action can occur.  This avoids repeating
the same Hurwitz estimate in the prime and non-prime cases.

\subsection{The numerical reduction}

\begin{lemma}\label{lem:tamereduction}
Let $q=p^h$ be odd and let $X/\F_{q^2}$ be a non-Hermitian maximal
curve of genus $g\ge2$. Suppose that a tame subgroup
$G\le\Aut(X)$ acts transitively on $X(\F_{q^2})$. Then
\[
q=p\le19.
\]
\end{lemma}

\begin{proof}
Put
\[
N=\#X(\F_{q^2})=q^2+1+2gq,
\qquad
s=|G_P|,
\quad P\in X(\F_{q^2}).
\]
By transitivity,
\[
|G|=sN.
\]
Since $p\nmid|G|$, the usual Hurwitz bound applies:
\begin{equation}\label{eq:tamehurwitz-general}
sN\le84(g-1);
\end{equation}
see \cite[Theorem~11.56]{HKT}.

Since $X$ is not Hermitian, Lemma~\ref{lem:gaps}(ii) gives
\[
g\le\frac{(q-1)^2}{4}.
\]
If $q=3$, then $g\le1$, a contradiction. Thus we may assume
$q\ge5$.

For fixed $q$, the function
\[
t\longmapsto
\frac{q^2+1+2qt}{t-1}
=
2q+\frac{(q+1)^2}{t-1},
\qquad t>1,
\]
is decreasing. Hence
\begin{equation}\label{eq:ratio-general}
\frac{N}{g-1}
\ge
2q+4+\frac{16}{q-3}.
\end{equation}
For $q\ge41$ the right-hand side is greater than $84$, contradicting
\eqref{eq:tamehurwitz-general}.

Suppose $23\le q\le37$. Then $N/(g-1)>48$, so
\eqref{eq:tamehurwitz-general} gives $s=1$. By
Remark~\ref{rem:threebranch}, $X/G$ is rational with three branch
points. If $(a,b,c)$ is the signature, with
$2\le a\le b\le c$, then
\[
0<
1-\frac1a-\frac1b-\frac1c
=
\frac{2(g-1)}{|G|}
<
\frac1{24}.
\]
The only hyperbolic triple satisfying this inequality is $(2,3,7)$.
Here the strict inequality is essential: $(2,3,8)$ has defect exactly
$1/24$ and therefore does not occur. Thus
\[
|G|=84(g-1).
\]
For $q=27$ this signature is not tame in characteristic $3$.
For
\[
q\in\{23,25,29,31,37\},
\]
the equality $|G|=N$ gives
\[
g=\frac{q^2+85}{84-2q}.
\]
The corresponding values are
\[
\frac{307}{19},\quad
\frac{355}{17},\quad
\frac{463}{13},\quad
\frac{523}{11},\quad
\frac{727}{5},
\]
respectively, none of which is an integer. Hence no value
$23\le q\le37$ can occur.

The only odd non-prime prime power below $23$ is $q=9$. In this case
\[
\frac{N}{g-1}\ge\frac{74}{3},
\]
so \eqref{eq:tamehurwitz-general} gives $s\le3$. Since
$3\nmid|G|$, one has $s\in\{1,2\}$.

Suppose first that $s=1$. Then
\[
\frac{|G|}{g-1}
=
\frac{N}{g-1}
\ge
\frac{74}{3},
\]
and hence
\[
0<
D:=\frac{2(g-1)}{|G|}
\le
\frac{3}{37}
<
\frac1{12}.
\]
By Remark~\ref{rem:threebranch}, $X/G$ is rational with three branch
points. Let their ramification indices be
\[
2\le a\le b\le c.
\]
Since the action is tame in characteristic $3$, one has
$3\nmid abc$, and
\[
D=
1-\frac1a-\frac1b-\frac1c.
\]
Thus
\[
\frac1a+\frac1b+\frac1c
\ge
\frac{34}{37}.
\]
If $a\ge4$, then the left-hand side is at most $3/4$, a
contradiction. Hence $a\le3$; since $3\nmid a$, we get $a=2$. If $b\ge5$, then
\[
\frac1a+\frac1b+\frac1c
\le
\frac12+\frac25
=
\frac9{10}
<
\frac{34}{37},
\]
again a contradiction. Since $3\nmid b$, we have $b=2$ or $4$.
The case $b=2$ would give
\[
\frac1a+\frac1b+\frac1c>1,
\]
which is impossible. Hence $b=4$. Finally,
\[
\frac1c
\ge
\frac{34}{37}-\frac34
=
\frac{25}{148},
\]
so $c\le5$. Since $c\ge4$ and $3\nmid c$, we have
$c=4$ or $5$. The value $c=4$ gives $D=0$, impossible since
$g\ge2$. Therefore the signature is $(2,4,5)$.

Consequently
\[
|G|=40(g-1).
\]
Since $s=1$,
\[
40(g-1)=N=82+18g,
\]
which gives
\[
g=\frac{61}{11},
\]
a contradiction.

If $s=2$, then
\[
\frac{|G|}{g-1}
=
\frac{2N}{g-1}
\ge
\frac{148}{3}
>
48.
\]
As above, the only possible signature is $(2,3,7)$, which is not
tame in characteristic $3$. Thus $q=9$ is impossible.

Hence $q$ is prime and $q\le19$.
\end{proof}



\subsection{The prime-field candidates}

Assume from now on that $q=p\le19$. The case $p=3$ is already
impossible for a non-Hermitian curve of genus at least two, since
Lemma~\ref{lem:gaps}(ii) gives $g\le1$. Thus the preceding lemma leaves
only
\[
p\in\{5,7,11,13,17,19\}.
\]

\begin{lemma}\label{lem:tamecandidates}
Let $X/\F_{p^2}$ be a non-Hermitian maximal curve of genus $g\ge2$,
and suppose that a tame subgroup $G\le\Aut(X)$ acts transitively on
$X(\F_{p^2})$. Put
\[
N=p^2+1+2gp,\qquad
s=|G_P|.
\]
Then the possibilities are exactly the seven rows in
Table~\ref{tab:candidates}.
\end{lemma}

\begin{table}[ht]
\centering
\caption{Numerical possibilities in the tame prime-field case.}
\label{tab:candidates}
\begin{tabular}{ccccc}
\toprule
$p$ & $g$ & $N$ & $s$ & signature\\
\midrule
5  & 3  & 56   & 3 & $(2,3,7)$\\
7  & 5  & 120  & 1 & $(2,3,10)$\\
11 & 9  & 320  & 1 & $(2,4,5)$\\
11 & 19 & 540  & 1 & $(2,3,10)$\\
11 & 25 & 672  & 3 & $(2,3,7)$\\
13 & 15 & 560  & 1 & $(2,4,5)$\\
19 & 41 & 1920 & 1 & $(2,3,8)$\\
\bottomrule
\end{tabular}
\end{table}

\begin{proof}
Let
\[
R:=\frac{|G|}{g-1}=\frac{sN}{g-1}.
\]
From \eqref{eq:ratio-general} with $q=p$ we have
\[
\frac{N}{g-1}\ge 2p+4+\frac{16}{p-3},
\]
and hence $R\ge22$. By Remark~\ref{rem:threebranch}, $X/G$ is
rational and the cover has exactly three branch points. Let
$(a,b,c)$ be its signature, with $2\le a\le b\le c$. Then
\[
R=\frac{2}{1-\frac1a-\frac1b-\frac1c},
\]
so
\begin{equation}\label{eq:signature-bound}
0<1-\frac1a-\frac1b-\frac1c\le\frac1{11}.
\end{equation}
The elementary triangle calculation gives
\[
(2,3,c),\qquad 7\le c\le13,
\]
together with
\[
(2,4,5),\qquad (2,4,6),\qquad (3,3,4).
\]
For these signatures,
\[
\begin{array}{c|c}
\text{signature} & R\\ \hline
(2,3,c) & \dfrac{12c}{c-6}\\[2mm]
(2,4,5) & 40\\
(2,4,6) & 24\\
(3,3,4) & 24.
\end{array}
\]

Combining \eqref{eq:tamehurwitz-general} with
\eqref{eq:ratio-general} gives
\[
\begin{array}{c|cccccc}
p&5&7&11&13&17&19\\ \hline
s\le&3&3&3&2&2&1.
\end{array}
\]
Finally,
\[
R(g-1)=s(p^2+1+2pg)
\]
gives
\begin{equation}\label{eq:g-from-R}
g=\frac{R+s(p^2+1)}{R-2sp}.
\end{equation}

If $s>1$, then the orbit $X(\F_{p^2})$ is a short orbit of $G$.
Hence its image in $X/G$ is a branch point, and therefore
\[
s\in\{a,b,c\}.
\]

The remaining step is a direct substitution into the ten signatures
listed above and the indicated finite ranges for $s$; no further
group-theoretic input is used. For each such pair we impose the
necessary conditions:
\[
p\nmid abc,\qquad
R-2sp>0,\qquad
g=\frac{R+s(p^2+1)}{R-2sp}\in\mathbf Z,\quad g\ge2,
\]
together with $s\in\{a,b,c\}$ when $s>1$.

For reference, before applying the genus-spectrum restriction the
integral solutions are the following; an entry $(g,s;a,b,c)$ records the
genus, stabilizer order and signature:
\[
\begin{array}{c|l}
5  &(3,3;2,3,7),\ (7,3;2,3,8),\ (19,3;2,3,9)\\
   &(19,2;2,3,12),\ (19,2;2,4,6)\\[1mm]
7  &(5,1;2,3,10),\ (17,2;2,3,9),\ (33,3;2,3,8)\\
   &(65,2;2,3,10)\\[1mm]
11 &(9,1;2,4,5),\ (19,1;2,3,10),\ (25,3;2,3,7)\\
   &(73,2;2,3,8),\ (73,1;2,3,12),\ (73,1;2,4,6)\\
   &(73,1;3,3,4),\ (505,1;2,3,13)\\[1mm]
13 &(15,1;2,4,5),\ (50,1;2,3,10),\ (491,1;2,3,11)\\
17 &(55,1;2,4,5),\ (163,1;2,3,9)\\
19 &(41,1;2,3,8),\ (201,1;2,4,5).
\end{array}
\]

The upper genus spectrum gives one further essential restriction. Put
\[
g_1=\frac{(p-1)^2}{4},
\qquad
g_2=\left\lfloor\frac{p^2-p+4}{6}\right\rfloor.
\]
By Lemma~\ref{lem:gaps}, an admissible non-Hermitian genus satisfies
\[
g=g_1\qquad\text{or}\qquad g\le g_2.
\]
Substitution in \eqref{eq:g-from-R}, followed by the genus-spectrum
restriction, gives exactly the seven rows of Table~\ref{tab:candidates}.
For $p=5$ one has $g_1=g_2=4$, so the distinction between the two
alternatives in the genus-spectrum restriction is vacuous in that
case. Notice that the genus-spectrum condition is essential for the
larger primes; for example, the formal solution
$(p,g)=(17,55)$ is excluded since
\[
46<55<64,
\]
where $g_2=46$ and $g_1=64$. The row $(p,g)=(11,19)$ lies exactly on
the third genus bound $g_2=19$.
\end{proof}

One row is removed by the second-genus classification.

\begin{lemma}\label{lem:gapeliminate}
The row $(p,g)=(11,25)$ in Table~\ref{tab:candidates} does not occur.
\end{lemma}

\begin{proof}
Here
\[
g=\frac{(p-1)^2}{4}.
\]
By Lemma~\ref{lem:gaps}(iii),
\[
X\simeq \cH_6:\quad y^6=x^{11}+x.
\]
By \eqref{eq:Hmaut},
\[
|\Aut(X)|=6\cdot11(11^2-1)=7920.
\]
Transitivity in this row would give
\[
|G|=3\cdot672=2016.
\]
But $7\mid2016$, whereas $7\nmid7920$, so
$|G|\nmid|\Aut(X)|$, a contradiction.
\end{proof}

The next three rows are excluded by a combination of the
characteristic-zero order bounds and a small group-theoretic argument.

\begin{lemma}\label{lem:triangleexclude}
The rows
\[
(p,g)=(11,19),\qquad (13,15),\qquad (19,41)
\]
do not occur.
\end{lemma}

\begin{proof}
In all three rows $s=1$, so the relevant data are
\[
\begin{array}{c|c|c}
g&|G|&\text{signature}\\ \hline
19&540 &(2,3,10)\\
15&560 &(2,4,5)\\
41&1920&(2,3,8).
\end{array}
\]
By Subsection~\ref{subsec:tame-specialization}, each action gives a
characteristic-zero action with the same genus, group order and signature.

For two rows no group-by-group classification is needed. Breuer's summary
table \cite[p.~91]{Breuer} gives the maximal orders of groups of
holomorphic automorphisms. The same values are recorded in Conder's
orientation-preserving table \cite[Case~A, genera 15 and 41]{ConderMax}:
\[
|\Aut|_{\max}=504\quad\text{in genus }15\quad\text{(type $(2,3,9)$)},
\]
\[
|\Aut|_{\max}=960\quad\text{in genus }41\quad\text{(type $(2,4,6)$)}.
\]
The orientation-preserving qualification matters here: Conder's second
list allows orientation-reversing automorphisms and gives order $1920$
in genus $41$, whereas automorphisms of the complex curve are
holomorphic and hence orientation-preserving. Thus the actions of orders
$560$ and $1920$ are impossible.

It remains to exclude the genus-$19$ row. We do this without using the
low-genus classification. Suppose that a group $G$ of order $540$ admits
a generating pair $x,y$ with
\begin{equation}\label{eq:2310-generators}
|x|=2,\qquad |y|=3,\qquad |xy|=10.
\end{equation}
Such a pair is supplied by the signature $(2,3,10)$; equivalently,
$G$ is a quotient of the triangle group $\Delta(2,3,10)$.

We first claim that $G$ is not perfect. Otherwise choose a maximal normal
subgroup $M\triangleleft G$. Then $S=G/M$ is a nonabelian simple group.
By Burnside's $p^a q^b$ theorem, $|S|$ is divisible by all three primes
$2,3,5$. Moreover a Sylow $2$-subgroup of $S$ cannot have order $2$:
otherwise it is central in its normalizer and Burnside's normal
$2$-complement theorem contradicts simplicity. Since $|S|\mid540$, this
leaves
\[
|S|\in\{60,180,540\}.
\]

The cases $|S|=180$ and $540$ are impossible. Let
$P\in\Syl_5(S)$. Then
\[
n_5(S)\in\{6,36\}.
\]
If $n_5(S)=6$, the conjugation action on the six Sylow $5$-subgroups
is faithful. Since a nonabelian simple group has no quotient of order
$2$, its image lies in $A_6$. For $|S|=540$ this is impossible by
order, while for $|S|=180$ it would be an index-$2$ subgroup of $A_6$,
hence normal, contradicting the simplicity of $A_6$.

If $n_5(S)=36$, then $|N_S(P)|$ is $5$ or $15$. In either case
\[
P\le Z(N_S(P)).
\]
For $|N_S(P)|=15$ this also follows from
\[
N_S(P)/C_S(P)\hookrightarrow\Aut(P)\cong C_4.
\]
Burnside's normal $5$-complement theorem again contradicts simplicity.

Hence $|S|=60$, so $|M|=9$. Since $xy$ has order $10$,
\[
\langle xy\rangle\cap M=1,
\]
because the order of this intersection divides both $10$ and $9$.
Thus the image of $xy$ in $S$ still has order $10$. But a simple group
of order $60$ has no element of order $10$. Indeed, if $c$ had order
$10$ and $P=\langle c^2\rangle$, then $n_5(S)=6$, hence
$|N_S(P)|=10$; since $\langle c\rangle\le N_S(P)$, equality holds and
$P\le Z(N_S(P))$, again contradicting Burnside's normal
$5$-complement theorem. This proves that $G$ is not perfect.

The abelianization of the triangle group $\Delta(2,3,10)$ is $C_2$:
in the abelianization the relation $(xy)^{10}=1$, together with
$x^2=y^3=1$, forces $y=1$. Since $G$ is a nonperfect quotient of
$\Delta(2,3,10)$, it follows that
\[
G/G'\cong C_2,
\qquad
|G'|=270.
\]

The subgroup
\[
H=\langle y,xyx\rangle
\]
is normalized by $x$. It is also normalized by $y$, since $y\in H$;
therefore $H\triangleleft G$ because $G=\langle x,y\rangle$.
Moreover $H\le G'$, while $G/H$ is generated by the image of $x$
and hence has order at most $2$. Since $[G:G']=2$, it follows that
\[
H=G'.
\]
In particular, $G'$ is generated by the two elements $y$ and $xyx$,
both of order $3$.

Let $T\in\Syl_2(G')$. Since $|G'|=270$, one has $|T|=2$, and therefore
\[
T\le Z(N_{G'}(T)).
\]
Burnside's normal $2$-complement theorem gives a normal subgroup
$K\triangleleft G'$ of index $2$. Thus $G'$ has a nontrivial quotient
$C_2$. This is impossible because $G'$ is generated by the two elements
$y$ and $xyx$, both of order $3$, and every homomorphism from $G'$ to
$C_2$ kills both generators. The genus-$19$ row is therefore impossible.
\end{proof}

It remains to treat
\[
(p,g)=(7,5),\qquad (11,9),
\]
and then to identify the surviving case $(p,g)=(5,3)$.


\subsection{\texorpdfstring{The genus-five case in characteristic $7$}
{The genus-five case in characteristic 7}}

\begin{proposition}\label{prop:seven}
There is no $\F_{49}$-maximal curve of genus $5$ admitting a tame
subgroup of automorphisms which is transitive on its rational points.
\end{proposition}

\begin{proof}
By Table~\ref{tab:candidates},
\[
|G|=120,
\qquad
\text{signature }(2,3,10).
\]
By Subsection~\ref{subsec:tame-specialization}, the signature
$(2,3,10)$ gives generators $x,y\in G$ such that
\[
|x|=2,\qquad |y|=3,\qquad |xy|=10.
\]
Put $c=xy$. We first extract directly from these data the central
involution needed below.

Let $P\in\Syl_5(G)$. We claim that $n_5(G)=6$. Indeed, if $P$ were
normal, then in $G/P$, which has order $24$, the image of $c$ would
have order dividing both $10$ and $24$, hence order at most $2$.
Thus $G/P$ would be a quotient of the spherical triangle group
$\Delta(2,3,2)\simeq S_3$, impossible because $|G/P|=24$.
Hence $n_5(G)=6$.

Let
\[
\varphi:G\longrightarrow \operatorname{Sym}(\Syl_5(G))\simeq S_6
\]
be the conjugation action, and let $K=\ker(\varphi)$. Since
$|N_G(P)|=120/6=20$, one has $K\le N_G(P)$. Moreover $5\nmid|K|$:
otherwise, by normality of $K$, it would contain all six Sylow
$5$-subgroups of $G$, and hence at least $1+6(5-1)=25$ elements,
contrary to $|K|\le20$. Therefore
\[
|K|\mid4.
\]
Now $S_6$ has no element of order $10$. Since $K$ is a $2$-group, the
$5$-part of $c$ survives in the quotient, so $\varphi(c)$ has order
$5$. Consequently
\[
z:=c^5\in K,
\]
and $|K|$ is $2$ or $4$.

Suppose that $|K|=4$. Then $G/K$ has order $30$ and acts faithfully and
transitively on the six Sylow $5$-subgroups. Since $K$ is a $2$-group,
the Sylow $5$-subgroups of $G$ map bijectively to those of $G/K$; hence
the stabilizer of $\overline P=PK/K$ is
$N_{G/K}(\overline P)=N_G(P)/K$, of order $5$. Thus $\overline P$ is
self-normalizing. Burnside's normal $5$-complement theorem then
gives a normal subgroup $H$ of order $6$. The conjugation action of
$\overline P$ on $H$ is trivial, because a group of order $6$ is
isomorphic to either $C_6$ or $S_3$, whose automorphism group has order
$2$ or $6$, respectively. Hence $\overline P$ is normal in $G/K$, a
contradiction. Therefore
\[
K=\langle z\rangle\simeq C_2.
\]
In particular, $z=c^5$ is a central involution of $G$. Notice also that
$G/K$ is generated by the images of $x$ and $y$ with product of order
$5$, so it is a quotient of $\Delta(2,3,5)\simeq A_5$; since
$|G/K|=60$, one has $G/K\simeq A_5$. We shall only need the centrality
of $z$.

Let $Q\in X$ be a point above the branch point of index $10$, so that
\[
G_Q=\langle c\rangle.
\]
The fibre of this branch point consists of
\[
\frac{|G|}{|G_Q|}
=
\frac{120}{10}
=
12
\]
points. Since $z$ is central and $z\in G_Q$, it belongs to every
conjugate of $G_Q$. Consequently, $z$ fixes all twelve points in this
fibre.

Put
\[
Y=X/\langle z\rangle
\]
and let $g_Y$ be the genus of $Y$. Since the characteristic is odd,
the quotient by $z$ is tame. If $r$ denotes the number of fixed points
of $z$, Riemann--Hurwitz gives
\[
2g(X)-2
=
2(2g_Y-2)+r.
\]
As $g(X)=5$ and $r\ge12$, we obtain
\[
8=2(2g_Y-2)+r,
\]
or equivalently
\[
r=12-4g_Y.
\]
Hence $r\le12$. Therefore
\[
r=12,
\qquad
g_Y=0.
\]
Thus $X$ is hyperelliptic and $z$ is its hyperelliptic involution.

The hyperelliptic involution is unique, hence it is defined over
$\F_{49}$. Therefore the quotient map
\[
X\longrightarrow Y
\]
is defined over $\F_{49}$. Since $Y$ has genus zero and the image of any point of
$X(\F_{49})$ gives an $\F_{49}$-rational point of $Y$, one has
\[
Y\simeq\PP^1_{\F_{49}}.
\]
Each $\F_{49}$-rational point of $Y$ has at most two
$\F_{49}$-rational points above it, and consequently
\[
\#X(\F_{49})
\le
2\#\PP^1(\F_{49})
=
2(49+1)
=
100.
\]
On the other hand, maximality gives
\[
\#X(\F_{49})
=
49+1+2\cdot5\cdot7
=
120,
\]
a contradiction.
\end{proof}

\subsection{\texorpdfstring{The genus-nine case in characteristic $11$}
{The genus-nine case in characteristic 11}}

\begin{proposition}\label{prop:eleven}
There is no $\F_{121}$-maximal curve of genus $9$ admitting a tame
subgroup of automorphisms which is transitive on its rational points.
\end{proposition}

\begin{proof}
By Table~\ref{tab:candidates},
\[
|G|=320,\qquad \text{signature }(2,4,5).
\]
Because the signature is $(2,4,5)$, there are generators $x,y$ of
$G$ such that
\[
 |x|=2,\qquad |y|=4,\qquad |xy|=5.
\]
Put $c=xy$ and $S=\langle c\rangle$. We first extract from these data
exactly the group structure needed below.

The Sylow $5$-subgroup $S$ is not normal. Indeed, if $S\triangleleft G$,
then in $G/S$ one has $\bar x\bar y=1$, so $G/S$ is generated by the
single involution $\bar x$; this contradicts $|G/S|=64$. Hence
$O_5(G)=1$. Since $|G|=2^6\cdot5$, Burnside's $p^a q^b$ theorem shows
that $G$ is solvable. Let
\[
        F=F(G)=O_2(G)
\]
be its Fitting subgroup. Then $F\ne1$ and, by the standard
self-centralizing property of the Fitting subgroup of a finite solvable
group,
\[
        C_G(F)\le F;
\]
see, for example, \cite{Isaacs}.

The abelianization of the triangle group $\Delta(2,4,5)$ is $C_2$.
Indeed, in the abelianization the relation $(xy)^5=1$, together with
$x^2=y^4=1$, gives $\bar x+\bar y=0$. Since $G$ is a nontrivial
solvable quotient of $\Delta(2,4,5)$, it is not perfect, and therefore
\[
        G/G'\simeq C_2.
\]

The subgroup $S$ acts nontrivially on $F$, for otherwise
$S\le C_G(F)\le F$. Let
\[
 d=\dim_{\F_2} F/\Phi(F).
\]
By the Burnside basis theorem, an automorphism of odd order acting
trivially on $F/\Phi(F)$ acts trivially on $F$. Hence the induced
$S$-action on $F/\Phi(F)$ is nontrivial, and therefore faithful. Thus
$5\mid |GL(d,2)|$, which forces $d\ge4$. Consequently
$|F|\ge16$, so
\[
        |F|\in\{16,32,64\}.
\]

The cases $|F|=16$ and $64$ are impossible. If $|F|=64$, then
$G/F$ has order $5$; the images of $x$ and $y$ are both trivial, so
they cannot generate $G/F$. If $|F|=16$, put $H=G/F$. Then
$|H|=20$, and the Sylow $5$-subgroup $\bar S$ of $H$ is normal. Since
$H/\bar S$ has order $4$, it is abelian, so $H'\le\bar S$ and
$|H'|\le5$. On the other hand, $H/H'$ is a quotient of $G/G'\simeq
C_2$, which would force $|H'|\ge10$, a contradiction. Therefore
\[
        |F|=32.
\]
It follows that $G/F$ has order $10$. Its abelianization is a quotient
of $C_2$, so it is nonabelian; hence
\[
        G/F\simeq D_5,
\]
where $D_5$ denotes the dihedral group of order $10$.

It remains to locate a central involution. Since $d\ge4$ and
$|F|=32$, either $d=4$ or $d=5$. If $d=4$, then $|\Phi(F)|=2$.
Thus
\[
        Z:=\Phi(F)
\]
is characteristic in $F$, hence normal in $G$; being of order $2$, it
is central in $G$. Moreover $F/Z\simeq C_2^4$.

Suppose now that $d=5$. Then $F\simeq C_2^5$. Since in this case
$C_G(F)=F$, the quotient $G/F\simeq D_5$ acts faithfully on the
$5$-dimensional $\F_2$-space $F$. A nontrivial irreducible
$\F_2C_5$-module has dimension
\[
        \operatorname{ord}_5(2)=4.
\]
By Maschke's theorem, the restriction to the normal subgroup
$C_5\triangleleft D_5$ therefore decomposes as a nontrivial
$4$-dimensional irreducible module plus a $1$-dimensional fixed space.
Let $Z$ be this fixed space. Then $|Z|=2$, and because $C_5$ is normal
in $D_5$, the subspace $Z$ is $D_5$-invariant. Hence again
$Z\triangleleft G$, so $Z\le Z(G)$, and $F/Z\simeq C_2^4$.

Thus in both cases, with
\[
        \overline G=G/Z,\qquad V=F/Z,
\]
we have
\[
        Z\simeq C_2,\qquad V\simeq C_2^4,\qquad
        \overline G/V\simeq D_5.
\]
Moreover a Sylow $5$-subgroup acts faithfully on $V$: for $d=4$ this
is the faithful $S$-action on $F/\Phi(F)$ above, while for $d=5$ it is the
nontrivial irreducible $4$-dimensional summand. No splitting assertion
for
\[
1\longrightarrow V\longrightarrow \overline G
\longrightarrow D_5\longrightarrow1
\]
is needed below. Let $z$ be the nontrivial element of $Z$.

We first show that $z$ is fixed-point-free. Since the action is tame,
every point stabilizer is cyclic. If $z$ fixed a point $P$, then,
because $z$ is central, it would fix every point of the orbit
$G\cdot P$. The only possible nontrivial stabilizer orders are
$2,4,5$. Since $z$ has order $2$, a stabilizer containing $z$ has
order $2$ or $4$, and therefore
\[
|G\cdot P|\ge\frac{320}{4}=80.
\]
On the other hand, if $r_z$ denotes the number of fixed points of $z$,
Riemann--Hurwitz for the tame double cover
$X\to X/\langle z\rangle$ gives
\[
r_z\le2g(X)+2=20.
\]
This is impossible. Hence $z$ is fixed-point-free.

Set
\[
Y=X/\langle z\rangle.
\]
The covering $X\to Y$ is unramified of degree $2$, so
Riemann--Hurwitz gives
\[
2\cdot9-2=2(2g(Y)-2),
\]
and hence
\[
g(Y)=5.
\]
By \cite[Theorem~3.10]{GSY}, the involution $z$ is defined over
$\F_{121}$. Hence Lemma~\ref{lem:quotientmax} applies, and $Y$ is
maximal over $\F_{121}$. In particular, $Y$ is supersingular.

Recall that
\[
\overline G=G/\langle z\rangle.
\]
The induced action of $\overline G$ on $Y$ is faithful: an element of
$G$ acting trivially on $Y=X/\langle z\rangle$ belongs to the deck group
$\langle z\rangle$.
Let $\overline P$ be the image of $P\in X$ in $Y$. If
$\overline g\in\overline G$ fixes $\overline P$, then a representative
$g\in G$ satisfies
\[
g(P)=P\qquad\text{or}\qquad g(P)=z(P).
\]
In the second case $zg$ fixes $P$. Thus every element of
$\overline G_{\overline P}$ has a representative in $G_P$. Since
$z$ is fixed-point-free,
\[
G_P\cap\langle z\rangle=1,
\]
and consequently the natural map induces an isomorphism
\[
G_P\xrightarrow{\;\sim\;}\overline G_{\overline P}.
\]
Therefore passage to $\overline G$ preserves all point-stabilizer
orders, and the signature of the $\overline G$-action on $Y$ remains
$(2,4,5)$.

Let $V\simeq C_2^4$ be the normal subgroup identified above and put
\[
Y_0=Y/V.
\]
A point stabilizer in $V$ is a subgroup of a point stabilizer in
$\overline G$, and hence is cyclic. Since $V$ is elementary abelian,
such a stabilizer is therefore either trivial or of order $2$. Let
\[
h=g(Y_0)
\]
and let $r$ be the number of branch points of the cover
$Y\to Y_0$. Each branch point has inertia group of order $2$, so
Riemann--Hurwitz gives
\[
2g(Y)-2
=
16(2h-2)+r\cdot16\left(1-\frac12\right).
\]
Since $g(Y)=5$, this becomes
\[
8=16(2h-2)+8r,
\]
and therefore
\[
r=5-4h.
\]
The only possibilities are
\[
(h,r)=(0,5)\qquad\text{or}\qquad(h,r)=(1,1).
\]
The latter is impossible. Indeed, for a connected tame abelian cover
of a genus-one curve with one branch point, the fundamental-group
relation has the form
\[
[a,b]\,c=1,
\]
where $c$ is an inertia generator. Since the Galois group is abelian,
this forces $c=1$, contrary to ramification. Hence
\[
Y/V\simeq\PP^1,\qquad r=5.
\]
Moreover the induced action of $\overline G/V$ on $Y/V$ is faithful,
since the kernel in $\overline G$ of the action on the quotient is
precisely the deck group $V$.

Let
\[
v_1,\ldots,v_5\in V
\]
be the nontrivial inertia generators at the five branch points. Since
the $V$-cover is connected, the $v_i$ generate $V$. The tame
fundamental-group relation gives
\[
v_1+\cdots+v_5=0.
\]
As $V$ has dimension $4$ over $\F_2$, the five generators have a
one-dimensional space of relations. Thus the displayed relation is
their unique nontrivial linear relation.

The quotient
\[
\overline G/V\simeq D_5
\]
acts on the five branch points. By the structural argument above, a Sylow
$5$-subgroup
\[
C_5\le\overline G
\]
acts faithfully on $V$. Since
\[
\operatorname{ord}_5(2)=4,
\]
the unique nontrivial irreducible $\F_2C_5$-module has dimension $4$.
Because $2\nmid5$, Maschke's theorem applies.  A faithful module must
contain the nontrivial irreducible $4$-dimensional summand; since $V$
itself has dimension $4$, there is no room for a trivial summand.
Hence $V$ is this nontrivial irreducible module. In particular,
\[
V^{C_5}=0.
\]

We claim that $C_5$ acts transitively on the five branch points. If it
fixed one of them, then it would normalize the corresponding inertia
subgroup $\langle v_i\rangle$. Since this subgroup has order $2$, its
unique nontrivial element $v_i$ would be fixed by $C_5$, contradicting
$V^{C_5}=0$. Hence $C_5$ has no fixed point on the set of five branch
points. An action of a group of order $5$ on a set of five elements
without fixed points is necessarily transitive. Consequently the
$D_5$-action is transitive as well.

We now work over $\overline{\F}_{11}$; no descent of the following
coordinate normalization is needed, since ordinarity and
supersingularity are geometric properties. Choose a generator $\rho$ of
the normal subgroup $C_5\triangleleft D_5$. Since $\rho$ is tame, it
has two fixed points on $\PP^1$; after a change of coordinate we may
write
\[
\rho(t)=\zeta t,
\]
where $\zeta$ is a primitive fifth root of unity. The five branch
points form a single $C_5$-orbit. Since $0$ and $\infty$ are the only
fixed points of $\rho$, neither belongs to this orbit. Thus, if
$\alpha$ is one branch point, the orbit is
\[
        \alpha\mu_5.
\]
After replacing $t$ by $t/\alpha$, we may therefore assume that the
five branch points are precisely
\[
\mu_5.
\]
Since $5\mid10=|\F_{11}^{*}|$, all fifth roots of unity lie in
$\F_{11}$; explicitly,
\[
\mu_5=\{1,3,4,5,9\}.
\]

Relabel the inertia generators so that $v_1$ corresponds to $t=1$.
Since
\[
v_1+\cdots+v_5=0
\]
is the unique relation among them, there exists a character
\[
\chi:V\longrightarrow C_2
\]
satisfying
\[
\chi(v_1)=0,\qquad
\chi(v_i)=1\quad(2\le i\le5).
\]
Indeed, the prescribed values respect the unique relation, since
\[
0+1+1+1+1=0
\qquad\text{in }\F_2.
\]
Put
\[
K=\ker(\chi).
\]
Then $E=Y/K$ is a double cover of $Y/V\simeq\PP^1$, branched exactly
at the four points
\[
3,4,5,9.
\]
Thus $E$ has genus $1$.

Over $\overline{\F}_{11}$, after the above change of coordinate and a
rescaling of the covering coordinate, $E$ has the model
\begin{equation}\label{eq:ellipticquotient}
E:\qquad
w^2=(t-3)(t-4)(t-5)(t-9)
=\frac{t^5-1}{t-1}
=t^4+t^3+t^2+t+1.
\end{equation}
For an odd-characteristic hyperelliptic model $w^2=f(t)$, the
Cartier operator on the regular differential $dt/w$ is nonzero exactly
when the coefficient of $t^{p-1}$ in $f(t)^{(p-1)/2}$ is nonzero; see,
for example, \cite[\S2]{Yui}. In genus one this is equivalent to
ordinarity. Thus, in characteristic
$11$, it suffices to compute
\[
[t^{10}]\,(t^4+t^3+t^2+t+1)^5
=
381
\equiv7\pmod{11}.
\]
Hence $E$ is ordinary.

On the other hand, $E$ is a quotient of the supersingular curve $Y$.
The induced maps on Jacobians give
\[
J(E)\xrightarrow{\pi^*}J(Y)
\xrightarrow{\operatorname{Nm}_\pi}J(E),
\]
with
\[
\operatorname{Nm}_\pi\circ\pi^*
=
[\deg(\pi)]
=
[|K|]
=
[8].
\]
Since $11\nmid8$, the image of $\pi^*$ is an abelian subvariety of
$J(Y)$ isogenous to $J(E)$. The Newton slopes of a supersingular
abelian variety, and hence of each of its abelian subvarieties up to
isogeny, are all $1/2$; see, for example, \cite[Chapter~I]{LiOort}.
Therefore $E$ is supersingular, contradicting its ordinarity.

Therefore no such curve $X$ exists.
\end{proof}


\subsection{The Klein quartic}

\begin{proposition}\label{prop:klein}
Every $\F_{25}$-maximal curve of genus $3$ is $\F_{25}$-isomorphic to
the $\F_{25}$-maximal form of Elkies' classical $S_4$-model of the
Klein quartic \cite[(1.11)]{ElkiesKlein}; see also
\cite[\S3.1]{BTT}. For this curve the full automorphism group acts
transitively on its $56$ rational points.
\end{proposition}

\begin{proof}
The explicit classification in \cite[Lemma~3.1(3)]{BMT} states that
an $\F_{25}$-maximal curve of genus $3$ is unique up to
$\F_{25}$-isomorphism. We use only this uniqueness statement here.
The terminology and the classical plane equation for the
``$S_4$-model'' go back to Elkies \cite[(1.11)]{ElkiesKlein}.  In his
notation the model is defined over
\[
        k=\mathbb{Q}(\sqrt{-7})=\mathbb{Q}(\alpha),
        \qquad
        \alpha=\frac{-1+\sqrt{-7}}2,
\]
see \cite[\S1.3]{ElkiesKlein}.  Since
$\alpha^2+\alpha+2=0$ and the discriminant $-7\equiv3\pmod 5$ is a
nonsquare in $\F_5$, reduction at $5$ gives coefficient field
$\F_{25}$.  Bootsma--Tafazolian--Top \cite[\S3.1]{BTT} exhibit this
$S_4$-model as an $\F_{25}$-maximal curve of genus $3$. Hence every
$\F_{25}$-maximal curve of genus $3$ is $\F_{25}$-isomorphic to this
model. The terminology ``$S_4$-model'' refers to this classical choice
of coordinates, in which a maximal subgroup $S_4<PSL(2,7)$ is especially
transparent, and not to the full automorphism group of the curve.

The same model carries a faithful action of $PSL(2,7)$; see
\cite[\S3.1]{BTT}. Hence $|\Aut(X)|\ge168$. On the other hand,
Roquette's theorem applies because $p=5>g+1=4$: a genus-$3$ curve in
characteristic $5$ satisfies the Hurwitz bound
\[
        |\Aut(X)|\le84(g-1)=168,
\]
unless it is the exceptional Roquette curve, which in characteristic
$5$ has genus $(5-1)/2=2$; see \cite{Roquette}. Consequently
\[
\Aut(X)\cong PSL(2,7),
\qquad
|\Aut(X)|=168.
\]
Moreover, since $X$ is maximal, every geometric
automorphism of $X$ is defined over $\F_{25}$ by
\cite[Theorem~3.10]{GSY}. In particular,
$X(\F_{25})$ is invariant under $\Aut(X)$.

Since $5\nmid168$, the action is tame. The Riemann--Hurwitz formula
gives
\[
\frac{2g(X)-2}{|\Aut(X)|}
=
\frac{4}{168}
=
\frac1{42}.
\]
It follows that $X/\Aut(X)$ is rational with three branch points and
the signature is
\[
(2,3,7).
\]
Consequently the three short orbits have lengths
\[
\frac{168}{2}=84,
\qquad
\frac{168}{3}=56,
\qquad
\frac{168}{7}=24,
\]
while every other orbit has length $168$.

Maximality gives
\[
\#X(\F_{25})
=
25+1+2\cdot3\cdot5
=
56.
\]
Since $X(\F_{25})$ is $\Aut(X)$-invariant, it is a disjoint union of
full $\Aut(X)$-orbits. The only way to express $56$ as a sum of orbit
lengths from
\[
\{24,56,84,168\}
\]
is the single orbit of length $56$. Hence $\Aut(X)$ acts transitively
on $X(\F_{25})$.
\end{proof}

\begin{remark}
Here the phrase ``the $\F_{25}$-maximal Klein quartic'' refers to this
$\F_{25}$-isomorphism class; it should not be confused with a particular
standard plane equation of the Klein quartic over a smaller field.
\end{remark}

\begin{proposition}\label{prop:tamefinal}
Let $q=p^h$ be odd and let $X/\F_{q^2}$ be a non-Hermitian maximal
curve of genus at least two. Suppose that a tame subgroup
$G\le\Aut(X)$ is transitive on $X(\F_{q^2})$. Then
\[
q=5,\qquad g=3,
\]
and $X$ is the $\F_{25}$-maximal Klein quartic.

Conversely, the $\F_{25}$-maximal Klein quartic satisfies these
conditions with $G=\Aut(X)\cong PSL(2,7)$.
\end{proposition}

\begin{proof}
Lemma~\ref{lem:tamereduction} gives
\[
q=p\le19,
\]
so in particular $h=1$. The case $p=3$ is impossible by
Lemma~\ref{lem:gaps}(ii), since a non-Hermitian maximal curve over
$\F_9$ would satisfy
\[
g\le\frac{(3-1)^2}{4}=1.
\]
Thus
\[
p\in\{5,7,11,13,17,19\}.
\]

Lemma~\ref{lem:tamecandidates} reduces these possibilities to the
seven rows of Table~\ref{tab:candidates}; in particular, no row with
$p=17$ remains. Lemma~\ref{lem:gapeliminate} excludes $(p,g)=(11,25)$,
while Lemma~\ref{lem:triangleexclude} excludes
\[
(p,g)=(11,19),\qquad (13,15),\qquad (19,41).
\]
Proposition~\ref{prop:seven} excludes $(p,g)=(7,5)$, and
Proposition~\ref{prop:eleven} excludes $(p,g)=(11,9)$. Hence the only
remaining possibility is
\[
(p,g)=(5,3).
\]
Proposition~\ref{prop:klein} then identifies $X$ with the
$\F_{25}$-maximal Klein quartic.

Conversely, Proposition~\ref{prop:klein} shows that the full
automorphism group of the $\F_{25}$-maximal Klein quartic is
$PSL(2,7)$ of order $168$ and acts transitively on its $56$ rational
points. Since $5\nmid168$, this action is tame.
\end{proof}


\section{The wild case}

Throughout this section, $q=p^h$ is odd, $X/\F_{q^2}$ is maximal of
genus $g\ge2$, and $G=\Aut(X)$ is the full geometric automorphism
group.  By \cite[Theorem~3.10]{GSY}, every element of $G$ is defined over
$\F_{q^2}$, so $G$ preserves
\[
          \Omega=X(\F_{q^2}).
\]
Assume that $G$ is transitive on $\Omega$ and that $p\mid|G|$.  Put
\[
          N=\#\Omega=q^2+1+2gq.
\]
Since $N\equiv1\pmod p$, the orbit--stabilizer identity
\[
          |G|=N|G_P|
\]
shows that $p\mid|G_P|$ for every $P\in\Omega$.

\subsection{\texorpdfstring{Fixed points and Sylow $p$-subgroups}{Fixed points and Sylow p-subgroups}}

By \cite[Lemma~1.1]{FGT}, the $\F_{q^2}$-Frobenius acts as $[-q]$ on
the Jacobian of a maximal curve; in particular, $X$ has $p$-rank zero.
We shall use the following consequence in the form proved in
\cite[Proposition~3.12]{GSY}.

\begin{lemma}\label{lem:pgroup-fixed}
Let $U\le\Aut(X)$ be a nontrivial $p$-subgroup.  Then $U$ fixes a unique
point $P\in\Omega$ and acts freely on $X\setminus\{P\}$.
\end{lemma}

\begin{proof}
By \cite[Proposition~3.12]{GSY}, $U$ fixes a unique point $P\in X$ and
acts freely on $X\setminus\{P\}$.  By \cite[Theorem~3.10]{GSY}, every
element of $U$ is defined over $\F_{q^2}$ and hence commutes with the
$\F_{q^2}$-Frobenius $F$.  Therefore $F(P)$ is again fixed by $U$.
By uniqueness, $F(P)=P$, so $P\in\Omega$.
\end{proof}

\begin{lemma}\label{lem:sylowbijection}
For $P\in\Omega$, put $S_P=G_P^{(1)}$.  Then $S_P$ is a Sylow
$p$-subgroup of $G$, and
\[
          P\longmapsto S_P
\]
is a $G$-equivariant bijection from $\Omega$ onto $\Syl_p(G)$.  Moreover,
\begin{equation}\label{eq:stabilizer}
          N_G(S_P)=G_P=S_P\rtimes C_m,
\end{equation}
where $C_m$ is cyclic, $(m,p)=1$, and $m$ is independent of $P$.  In
particular,
\begin{equation}\label{eq:number-sylow}
          |\Syl_p(G)|=N=q^2+1+2gq.
\end{equation}
Distinct Sylow $p$-subgroups have trivial intersection; equivalently,
$S_P$ is a TI subgroup of $G$.  Finally, $S_P$ acts freely on
$X\setminus\{P\}$.
\end{lemma}

\begin{proof}
Since $|G|=N|G_P|$ and $p\nmid N$, the $p$-parts of $|G|$ and $|G_P|$
coincide.  By \cite[Proposition~3.8.5]{Stichtenoth}, $G_P^{(1)}$ is a
normal $p$-subgroup of $G_P$ and $G_P/G_P^{(1)}$ is cyclic of order
prime to $p$.  Hence $S_P=G_P^{(1)}$ is the unique Sylow $p$-subgroup
of $G_P$, and therefore also a Sylow $p$-subgroup of $G$.  By
Schur--Zassenhaus,
\[
          G_P=S_P\rtimes C_m
\]
for a cyclic group $C_m$ of order prime to $p$.  Since $G$ is transitive
on $\Omega$, the stabilizers $G_P$ are conjugate; hence $m$ is
independent of $P$.

Conversely, let $S\in\Syl_p(G)$.  By Lemma~\ref{lem:pgroup-fixed}, $S$
has a unique fixed point $P\in\Omega$, so $S\le G_P$.  Both $S$ and
$S_P$ have order equal to the $p$-part of $|G_P|$; since $S_P$ is the
unique Sylow $p$-subgroup of $G_P$, it follows that $S=S_P$.  This proves
surjectivity.  Injectivity follows from the uniqueness of the fixed
point, and the map is $G$-equivariant because
\[
          S_{\sigma(P)}=\sigma S_P\sigma^{-1}
          \qquad(\sigma\in G).
\]
Thus \eqref{eq:number-sylow} follows immediately.

Since $S_P\triangleleft G_P$, we have $G_P\le N_G(S_P)$.  Conversely,
an element normalizing $S_P$ must preserve its unique fixed point, and
hence belongs to $G_P$.  Thus
\[
          N_G(S_P)=G_P.
\]

If $P\ne Q$ and $1\ne A\le S_P\cap S_Q$, then the nontrivial
$p$-subgroup $A$ fixes both $P$ and $Q$, contradicting
Lemma~\ref{lem:pgroup-fixed}.  Hence distinct Sylow $p$-subgroups have
trivial intersection, so $S_P$ is a TI subgroup of $G$.

Finally, Lemma~\ref{lem:pgroup-fixed} shows directly that $S_P$ acts
freely on all of $X\setminus\{P\}$.
\end{proof}

\subsection{The noncyclic Sylow case}

\begin{proposition}\label{prop:noncyclic}
If a Sylow $p$-subgroup of $G$ is noncyclic, then
\[
          X\simeq_{\F_{q^2}}\cH_{q+1}.
\]
\end{proposition}

\begin{proof}
Suppose, to the contrary, that $X$ is not Hermitian.  Then
Lemma~\ref{lem:gaps}(ii) gives
\begin{equation}\label{eq:nonherm-genus-wild}
          g\le\frac{(q-1)^2}{4}.
\end{equation}

Let $S$ be a noncyclic Sylow $p$-subgroup, with unique fixed point
$P\in\Omega$.  By Lemma~\ref{lem:sylowbijection},
\[
          N_G(S)=G_P=S\rtimes C_m,
\]
where $C_m$ is cyclic of order prime to $p$, and $S$ is a TI subgroup
of $G$.  Moreover,
\[
          [G:N_G(S)]=[G:G_P]=N>1,
\]
so $N_G(S)\ne G$.  Since $p$ is odd and $S$ is noncyclic, all hypotheses of
Theorem~\ref{thm:GMP-TI} are satisfied.  Hence, with
\[
          M:=Z(G_P),
\]
the subgroup $M$ is a normal $p'$-subgroup of $G$, the quotient $G/M$
is almost simple, and $G$ acts doubly transitively on $\Syl_p(G)$.

We first show that $M=1$.  Since $M\le G_P$, it fixes $P$; normality
then implies that $M$ fixes every point of the $G$-orbit $\Omega$.
If $M\ne1$, choose an element $\tau\in M$ of prime order
$\ell\ne p$.  If $f_\tau$ denotes the number of fixed points of
$\tau$, Riemann--Hurwitz for the tame cyclic cover
$X\to X/\langle\tau\rangle$ gives
\[
          f_\tau\le
          \frac{2g-2+2\ell}{\ell-1}
          \le 2g+2.
\]
On the other hand $\tau$ fixes all $N$ points of $\Omega$, while
\[
          N-(2g+2)=q^2-1+2g(q-1)>0,
\]
a contradiction.

Thus $G$ itself is almost simple.  The $G$-equivariant bijection in
Lemma~\ref{lem:sylowbijection} transfers the double transitivity on
$\Syl_p(G)$ to $\Omega$.  In particular,
\[
          |G|\ge N(N-1).
\]
By \eqref{eq:nonherm-genus-wild}, $q^2>4g$.  Since
$N>2gq$ and $N-1>2gq$,
\[
          |G|>4g^2q^2>16g^3>8g^3.
\]
Henn's classification \cite[Theorem~11.127]{HKT} now applies.  The
Suzuki and the characteristic-two hyperelliptic cases in that list are
excluded because $p$ is odd; the remaining odd-characteristic cases
are the geometrically Hermitian curves and the Roquette curves
\[
          y^2=x^r-x,\qquad r=p^a.
\]
A Roquette curve is hyperelliptic.  Its hyperelliptic involution is
central in the full automorphism group, by uniqueness of the
hyperelliptic $g^1_2$.  On the other hand an almost simple group is
centerless: if $T$ is its nonabelian simple socle, then
$Z(G)\le C_G(T)=1$.  Hence the Roquette case is impossible.  We are
left with
\[
          X_{\overline{\F}_p}\simeq\cH_{r+1}
\]
for some $r=p^a$.  Since $g=r(r-1)/2\ge2$, we have $r\ge3$.

Under this geometric identification $G\simeq PGU(3,r)$.  The latter
group acts transitively on the $r^3+1$ rational points of the Hermitian
curve; see \cite[Theorem~12.24(iv)]{HKT}.  A point stabilizer has a
unique Sylow $p$-subgroup by \cite[Lemma~11.44]{HKT}, and this is also
a Sylow $p$-subgroup of $G$ because its index $r^3+1$ is prime to $p$.
Since every nontrivial $p$-subgroup fixes a unique point, the Sylow
$p$-subgroups of $G$ are therefore in bijection with these points.
Thus $G$ has exactly $r^3+1$ Sylow $p$-subgroups.  By
Lemma~\ref{lem:sylowbijection},
\[
          r^3+1=N=q^2+1+r(r-1)q,
\]
and therefore
\begin{equation}\label{eq:r-q}
          r^3=q\bigl(q+r(r-1)\bigr).
\end{equation}

The maximal genus bound gives
\[
          \frac{r(r-1)}2=g\le\frac{q(q-1)}2,
\]
hence $r\le q$.  If $r<q$, write $r=p^a$ and $q=p^h$, with $a<h$.
Since $p\nmid r-1$,
\[
          v_p\bigl(q+r(r-1)\bigr)=a.
\]
Taking $p$-adic valuations in \eqref{eq:r-q} gives
$3a=h+a$, hence $h=2a$ and $q=r^2$.  Substitution into
\eqref{eq:r-q} gives
\[
          r^3=r^3(2r-1),
\]
a contradiction.  Therefore $r=q$.  Hence
$g=q(q-1)/2$, and Lemma~\ref{lem:gaps}(i) yields
\[
          X\simeq_{\F_{q^2}}\cH_{q+1},
\]
contrary to our assumption.  This proves the proposition.
\end{proof}

\begin{remark}\label{rem:ree-and-geometric-hermitian}
The final argument forcing $r=q$ is essential.  A curve may be
geometrically Hermitian with parameter $r<q$ and still be maximal over
$\F_{q^2}$.  For example,
\[
          \cH_4:\quad y^4=x^3+x
\]
has genus $3$ and is maximal over
$\F_{729}=\F_{27^2}$: its $\F_9$-Frobenius eigenvalues are all $-3$,
so it has
\[
          729+1+2\cdot3\cdot27=892
\]
$\F_{729}$-rational points.  Its full geometric automorphism group is
$PGU(3,3)$, which has only $3^3+1=28$ Sylow $3$-subgroups.  Thus it
cannot act transitively on those $892$ points.
\end{remark}

\subsection{The cyclic Sylow case}

The cyclic case requires a local ramification argument.  Unlike the
noncyclic case, there is no comparable group-theoretic classification
that by itself rules out cyclic TI Sylow $p$-subgroups; such groups occur
abundantly, for instance in $PSL(2,p)$.  We therefore use the lower
ramification filtration at the unique fixed point.

\begin{lemma}\label{lem:noncentral}
Assume that $X$ is non-Hermitian.  Let $S\in\Syl_p(G)$ be cyclic, let
$P$ be its unique fixed point, and write
\[
          |S|=n=p^a,\qquad
          G_P=S\rtimes C_m.
\]
Then the conjugation action of $C_m$ on $S$ is nontrivial.  If its image
has order $c$, then
\[
          1<c\mid p-1.
\]
\end{lemma}

\begin{proof}
Suppose, to the contrary, that the conjugation action of $C_m$ on $S$
is trivial.  Since $S$ is cyclic and $N_G(S)=G_P=S\rtimes C_m$, this
implies $S\le Z(N_G(S))$.  Burnside's normal $p$-complement theorem
\cite[Theorem~5.13]{Isaacs} gives a normal $p'$-subgroup $K\triangleleft G$ such that
$G=KS$ and $K\cap S=1$.  Since $S\le G_P$ and the quotient map
$G\to G/K\simeq S$ restricts surjectively to $G_P$, its kernel on $G_P$
is $K\cap G_P$.  Hence
\[
          [G_P:K\cap G_P]=n,
          \qquad |K\cap G_P|=m.
\]
Therefore
\[
          |K\cdot P|
          =\frac{|K|}{|K\cap G_P|}
          =\frac{|G|/n}{m}
          =N.
\]
Since $K\cdot P\subseteq G\cdot P=\Omega$, it follows that
$K\cdot P=\Omega$.  Thus the tame group $K$ is transitive on $\Omega$.
Proposition~\ref{prop:tamefinal} then forces $q=5$ and $X$ to be the
maximal Klein quartic.  But its full geometric automorphism group has
order $168$, which is prime to $5$, contradicting $p\mid|G|$.
Therefore the conjugation action is nontrivial.

Since $|\Aut(C_{p^a})|=p^{a-1}(p-1)$ and $C_m$ has order prime
to $p$, the order $c$ of the conjugation image satisfies
$1<c\mid p-1$.
\end{proof}

\begin{lemma}\label{lem:local}
Assume the notation of Lemma~\ref{lem:noncentral}, and let $c$ be the
order of the conjugation image of $C_m$ on $S$.  Let $j$ be the first
lower ramification jump of $S$, so that
\[
 G_P^{(1)}=\cdots=G_P^{(j)}=S,
 \qquad
 G_P^{(j+1)}<S.
\]
Then
\begin{equation}\label{eq:localrelation}
          \frac{m}{\gcd(m,j)}=c.
\end{equation}
Consequently
\begin{equation}\label{eq:Delta-lower}
          j\ge\frac{m}{c},
          \qquad
          \Delta\ge\frac{m(n-1)}{c},
\end{equation}
where
\[
          \Delta=\sum_{i\ge1}\bigl(|G_P^{(i)}|-1\bigr).
\]
\end{lemma}

\begin{proof}
Let $\tau$ generate $C_m$.  We use the standard structure of the
ramification filtration from
\cite[Proposition~3.8.5]{Stichtenoth}: $G_P^{(1)}$ is the unique
Sylow $p$-subgroup of $G_P$, while
$G_P^{(i)}/G_P^{(i+1)}$ is elementary abelian for $i\ge1$.

The action of $C_m$ on the cotangent line at $P$ is faithful.  Indeed,
an element of $C_m$ acting trivially on that line would belong to
$G_P^{(1)}=S$, and hence would be trivial.  Let $\zeta$ be the
eigenvalue of $\tau$ on the cotangent line; then $\zeta$ has order
$m$.  Starting from any local parameter $u$, set
\[
       t=\sum_{i=0}^{m-1}\zeta^{-i}\tau^i(u).
\]
Modulo the square of the maximal ideal one has
$t\equiv m u$, so $t$ is again a local parameter because $p\nmid m$.
Moreover,
\[
       \tau(t)=\zeta t.
\]

The quotient $S/G_P^{(j+1)}$ has order $p$, because it is both cyclic
and elementary abelian.  Choose
$\sigma\in S\setminus G_P^{(j+1)}$.  Then
\[
          \sigma(t)=t+a t^{j+1}+O(t^{j+2}),
          \qquad a\ne0.
\]
A direct calculation gives
\[
          \tau\sigma\tau^{-1}(t)
          =t+a\zeta^j t^{j+1}+O(t^{j+2}).
\]
The coefficient map identifies
$S/G_P^{(j+1)}$ with a one-dimensional $\F_p$-subspace of the additive
field, and conjugation by $\tau$ preserves this subspace.  Hence
multiplication by $\zeta^j$ restricts to an $\F_p$-linear automorphism
of it (in particular $\zeta^j\in\F_p^*$), so conjugation by $\tau$
acts on $S/G_P^{(j+1)}\simeq C_p$ by the scalar $\zeta^j$.

Since $S$ is cyclic and $S/G_P^{(j+1)}$ has order $p$,
$G_P^{(j+1)}=S^p$.  The reduction map
\[
          \Aut(C_{p^a})\longrightarrow\Aut(C_p)
\]
has $p$-group kernel.  Hence it is injective on the prime-to-$p$
conjugation image of $C_m$, whose order is $c$.  Therefore
\[
          \operatorname{ord}(\zeta^j)
          =\frac{m}{\gcd(m,j)}=c.
\]
The first $j$ positive ramification groups are equal to $S$, so
$\Delta\ge j(n-1)$, which gives \eqref{eq:Delta-lower}.

\end{proof}

\begin{proposition}\label{prop:cyclic}
There is no non-Hermitian maximal curve satisfying the transitivity
hypothesis for which a Sylow $p$-subgroup of $G$ is cyclic.
\end{proposition}

\begin{proof}
Assume that $X$ is non-Hermitian and that
\[
          S\in\Syl_p(G)
\]
is cyclic.  Let $P$ be its unique fixed point and write
\[
          |S|=n=p^a,\qquad
          G_P=S\rtimes C_m,\qquad
          s:=|G_P|=nm.
\]
By Lemma~\ref{lem:noncentral}, the conjugation image of $C_m$ on $S$
has order $c$ with
\[
          1<c\mid p-1.
\]
In particular $m\ge2$ and $s\ge2p$.

Since $X$ is not Hermitian, the same monotonicity calculation as in
\eqref{eq:ratio-general} gives
\[
 \frac{|G|}{g-1}
 =s\frac{N}{g-1}
 \ge
 2p\left(2q+4+\frac{16}{q-3}\right)>84.
\]
When $q=p=3$ the non-Hermitian genus bound already gives $g\le1$;
otherwise the displayed expression is defined and the inequality is
immediate.

We now apply the short-orbit classification
\cite[Theorem~2.4]{BMT}.  Since $|G|>84(g-1)$, the quotient $X/G$ is
rational and one of the following four configurations occurs: three
short orbits, one non-tame and two tame with stabilizer order $2$; two
non-tame short orbits; one non-tame short orbit; or two short orbits,
one non-tame and one tame. In particular, all Riemann--Hurwitz formulas
below are written with $g(X/G)=0$. We do not use the auxiliary bound
$|G|<8g^3$ occurring in the last case.

The only possible non-tame short orbit is $\Omega$.  Indeed, if
$p\mid|G_Q|$, choose a subgroup $A\le G_Q$ of order $p$ and a Sylow
$p$-subgroup $S'$ of $G$ containing $A$.  By
Lemma~\ref{lem:pgroup-fixed}, both $A$ and $S'$ have unique fixed
points, and the fixed point of $S'$ belongs to $\Omega$.  Since that
point is also fixed by $A$, it must be $Q$.  Hence $Q\in\Omega$.

Thus the case of two non-tame short orbits is impossible.  If
$\Omega$ were the only short orbit, then, using $|G|=sN$, the
Riemann--Hurwitz formula would give
\[
 \begin{aligned}
 2g-2
   &=-2|G|+N d_P\\
   &=N(d_P-2s),
 \end{aligned}
\]
which is impossible because the right-hand side is an integral
multiple of $N$, whereas $0<2g-2<N$.  If there were three short
orbits, case~\textup{(1)} of \cite[Theorem~2.4]{BMT} (recall that $p$
is odd) says that the other two are tame and that every point in them
has stabilizer of order $2$.  Thus each contributes $|G|/2$ to the
different, and hence
\[
 \begin{aligned}
 2g-2
   &=-2|G|+N d_P+|G|\\
   &=N(d_P-s),
 \end{aligned}
\]
again impossible because the right-hand side is an integral multiple
of $N$.  Therefore there are exactly two short orbits:
$\Omega$ and one tame short orbit.

Let $Q$ lie in the tame short orbit and put
\[
          e=|G_Q|.
\]
Then $p\nmid e$, $G_Q$ is cyclic, and the tame different exponent at
$Q$ is $d_Q=e-1$.  At $P$ put
\[
          \Delta=\sum_{i\ge1}\bigl(|G_P^{(i)}|-1\bigr).
\]
Hilbert's different formula gives
\[
          d_P=(|G_P|-1)+\Delta=(s-1)+\Delta.
\]
Since the tame orbit has length $|G|/e=sN/e$, Riemann--Hurwitz gives
\[
\begin{aligned}
 2g-2
   &=-2|G|+N d_P+\frac{|G|}{e}(e-1)\\
   &=-2sN+N\bigl((s-1)+\Delta\bigr)
      +\frac{sN}{e}(e-1)\\
   &=N\left(\Delta-1-\frac{s}{e}\right).
\end{aligned}
\]
Thus
\begin{equation}\label{eq:globalRH}
 \frac{2g-2}{N}
 =
 \Delta-1-\frac{s}{e}.
\end{equation}
The left-hand side belongs to $(0,1)$.  Therefore $e\nmid s$ and
\begin{equation}\label{eq:Delta}
          \Delta=\left\lceil\frac{s}{e}\right\rceil+1.
\end{equation}
Writing
\[
          s=\alpha e+\rho,
          \qquad 1\le\rho\le e-1,
\]
we obtain
\[
          e(2g-2)=N(e-\rho).
\]
In particular,
\begin{equation}\label{eq:e-lower}
          e\ge\frac{N}{2g-2}
          =q+\frac{(q+1)^2}{2g-2}>q.
\end{equation}
We shall also use
\begin{equation}\label{eq:ratio-less-1q}
          \frac{2g-2}{N}<\frac1q,
\end{equation}
which is equivalent to $q(2g-2)<N$ and follows immediately from
$N=q^2+1+2gq$.

Combining \eqref{eq:Delta} and \eqref{eq:Delta-lower}, and using
$s=nm$, we have
\[
 \frac{m(n-1)}{c}
 \le \Delta
 =\left\lceil\frac{nm}{e}\right\rceil+1
 <\frac{nm}{e}+2,
\]
where the last inequality is $\lceil x\rceil<x+1$ for nonintegral
$x$; here $e\nmid s$.  Hence
\begin{equation}\label{eq:master}
 m\left(\frac{n-1}{c}-\frac{n}{e}\right)<2.
\end{equation}

We distinguish the cases $q>p$ and $q=p$.

\medskip
\noindent\emph{Case 1: $q>p$.}
Then $q\ge p^2$.  Since $c\le p-1$, $e>q$, and $n\ge p$,
\[
 \frac{n-1}{c}-\frac{n}{e}
 >
 \frac{n-1}{p-1}-\frac{n}{q}.
\]
The function
\[
          f(x)=\frac{x-1}{p-1}-\frac{x}{q}
\]
is increasing because
$f'(x)=1/(p-1)-1/q>0$.  Hence, for $n\ge p$,
\[
 \frac{n-1}{p-1}-\frac{n}{q}
 \ge f(p)=1-\frac pq\ge1-\frac1p.
\]
Equation \eqref{eq:master} yields
\[
          m<\frac{2p}{p-1}.
\]
Since $c>1$ divides $m$, necessarily
\[
          m=c=2.
\]
Now \eqref{eq:Delta-lower} and \eqref{eq:Delta} imply
\[
          1-\frac3n<\frac2e<\frac2q.
\]
For $p\ge5$ this is impossible because $n\ge5$ and $q\ge25$.
For $p=3$ it is impossible if $n\ge9$.  Thus the only remaining
possibility is $p=3$, $n=3$.  Then $s=6$, and $e>q\ge9$ gives
$\Delta=2$.  Equation \eqref{eq:globalRH} becomes
\[
          \frac{2g-2}{N}=1-\frac6e>\frac13,
\]
whereas \eqref{eq:ratio-less-1q} gives
\[
          \frac{2g-2}{N}<\frac1q\le\frac19,
\]
a contradiction.

\medskip
\noindent\emph{Case 2: $q=p$.}
The non-Hermitian genus bound forces $p\ge5$.  From
$c\le p-1$ and $e>p$, equation \eqref{eq:master} gives
\[
          m\,\frac{n-p}{p(p-1)}<2.
\]
If $n\ge p^2$, then the left-hand side is at least $m\ge2$, a
contradiction.  Hence
\[
          n=p.
\]
There is now only one positive ramification jump, and
\[
          \Delta=j(p-1).
\]
Since $e>p$, the lower bound \eqref{eq:Delta-lower} and
\eqref{eq:Delta} give
\[
          m\le\Delta\le m+1.
\]

If $\Delta=m+1$, then
\[
          j(p-1)=m+1.
\]
Hence $\gcd(m,j)=1$, so \eqref{eq:localrelation} gives $c=m$.
Since $c=m$ and $c\mid p-1$, we have $m\le p-1$.  Thus
$j=1$ and $m=p-2$.  But then $p-2\mid p-1$, which forces $p=3$.
This is impossible because $p\ge5$.

Thus $\Delta=m$, so
\[
          m=j(p-1).
\]
Hence $\gcd(m,j)=j$, and \eqref{eq:localrelation} gives $c=p-1$.
Write $e=p+k$, $k\ge1$.  From \eqref{eq:globalRH} and \eqref{eq:ratio-less-1q},
\[
          0<
          \frac{2g-2}{N}
          =
          \frac{k(m-1)-p}{p+k}
          <\frac1p.
\]
Therefore
\begin{equation}\label{eq:kineq}
          k\bigl(p(m-1)-1\bigr)<p(p+1).
\end{equation}
For $p\ge7$, \eqref{eq:kineq} forces $k=1$: if $k\ge2$, then
$m\ge p-1$ gives
\[
  k\bigl(p(m-1)-1\bigr)
  \ge2\bigl(p(p-2)-1\bigr)>p(p+1).
\]
With $k=1$, positivity gives $m>p+1$.  Since $p-1\mid m$, it follows
that $m\ge2(p-1)$, whereas \eqref{eq:kineq} gives $m\le p+2$.
This is a contradiction.

It remains to consider $p=5$.  Equation \eqref{eq:kineq}, together
with $4\mid m$ and positivity, leaves only
\[
          m=4,\qquad k=2,\qquad e=7.
\]
Then
\[
          \frac{2g-2}{N}=\frac17.
\]
But a non-Hermitian maximal curve over $\F_{25}$ has $g\le4$.  Since
\[
          \frac{2g-2}{N}
          =\frac{2g-2}{26+10g}
\]
is increasing for $g\ge2$, it follows that
\[
          \frac{2g-2}{N}
          \le\frac{6}{66}
          =\frac1{11},
\]
a contradiction.

Thus the cyclic Sylow case is impossible for every odd $q$.
\end{proof}


\begin{proposition}\label{prop:wild}
Let $q=p^h$ be odd and let $X/\F_{q^2}$ be maximal of genus
$g\ge2$.  If $\Aut(X)$ is transitive on $X(\F_{q^2})$ and
$p\mid|\Aut(X)|$, then
$X\simeq_{\F_{q^2}}\cH_{q+1}$.
\end{proposition}

\begin{proof}
Let $G=\Aut(X)$ and let $S\in\Syl_p(G)$.  If $S$ is noncyclic,
Proposition~\ref{prop:noncyclic} gives
$X\simeq_{\F_{q^2}}\cH_{q+1}$.  If $S$ is cyclic and $X$ were
non-Hermitian, Proposition~\ref{prop:cyclic} would give a contradiction.
Hence $X\simeq_{\F_{q^2}}\cH_{q+1}$ also in this case.
\end{proof}

\section{Proof of the main theorem}

\begin{proof}[Proof of Theorem~\ref{thm:main}]
The Hermitian curve has automorphism group $PGU(3,q)$, which acts
doubly transitively on its $q^3+1$ rational points.  The $\F_{25}$-maximal Klein quartic is transitive by
Proposition~\ref{prop:klein}.

Conversely, let $G=\Aut(X)$ be transitive on $X(\F_{q^2})$.
If $p\mid|G|$, Proposition~\ref{prop:wild} gives the Hermitian curve.
Assume $p\nmid|G|$.  Then $X$ cannot be Hermitian, since
$PGU(3,q)$ has order divisible by $p$.  Proposition~\ref{prop:tamefinal}
therefore gives $q=5$, $g=3$, and the maximal Klein quartic.
\end{proof}

\section*{Acknowledgments}
OpenAI's ChatGPT was used to assist with language editing and the
preparation of the final version of the manuscript. All mathematical
arguments and references were checked by the author, who takes full
responsibility for the content.

\section*{Statements and Declarations}
\textbf{Funding.}
The author was partially supported by CNPq grant no.~302774/2025-4,
FAPESP grant no.~2024/00923-6, and FAEPEX grant no.~3485/25.

\textbf{Competing interests.}
The author has no relevant financial or non-financial interests to
disclose.

\begin{thebibliography}{99}

\bibitem{BMT}
D.~Bartoli, M.~Montanucci and F.~Torres,
\emph{$\F_{p^2}$-maximal curves with many automorphisms are
Galois-covered by the Hermitian curve},
Adv. Geom. \textbf{21} (2021), 325--336.

\bibitem{BTT}
S.~Bootsma, S.~Tafazolian and J.~Top,
\emph{Maximal curves and Tate--Shafarevich results for quartic and
sextic twists},
Finite Fields Appl. \textbf{91} (2023), 102256.

\bibitem{ElkiesKlein}
N.~D.~Elkies,
\emph{The Klein quartic in number theory},
in S.~Levy (ed.), \emph{The Eightfold Way: The Beauty of Klein's
Quartic Curve}, Math. Sci. Res. Inst. Publ. \textbf{35},
Cambridge University Press, Cambridge, 1999, 51--102.

\bibitem{Breuer}
T.~Breuer,
\emph{Characters and Automorphism Groups of Compact Riemann Surfaces},
London Math. Soc. Lecture Note Ser. \textbf{280},
Cambridge University Press, Cambridge, 2000.

\bibitem{ConderMax}
M.~D.~E.~Conder,
\emph{Maximum group orders by genus}, online data table,
\url{https://www.math.auckland.ac.nz/~conder/MaximumGroupOrdersByGenus-orientable.txt}.

\bibitem{FGT}
R.~Fuhrmann, A.~Garcia and F.~Torres,
\emph{On maximal curves},
J. Number Theory \textbf{67} (1997), no.~1, 29--51.

\bibitem{GK2010}
M.~Giulietti and G.~Korchm\'aros,
\emph{Automorphism groups of algebraic curves with $p$-rank zero},
J. Lond. Math. Soc. (2) \textbf{81} (2010), 277--296.

\bibitem{GK2019}
M.~Giulietti and G.~Korchm\'aros,
\emph{Algebraic curves with many automorphisms},
Adv. Math. \textbf{349} (2019), 162--211.

\bibitem{GKMSurvey}
M.~Giulietti, G.~Korchm\'aros and M.~Montanucci,
\emph{Maximal curves over finite fields, past, present and future},
Panoramas et Synth\`eses \textbf{60} (2023), 37--66.

\bibitem{GSY}
B.~Gunby, A.~Smith and A.~Yuan,
\emph{Irreducible canonical representations in positive characteristic},
Res. Number Theory \textbf{1} (2015), Art.~3, 25~pp.

\bibitem{GMP}
R.~Guralnick, B.~Malmskog and R.~Pries,
\emph{The automorphism groups of a family of maximal curves},
J. Algebra \textbf{361} (2012), 92--106.

\bibitem{HKT}
J.~W.~P.~Hirschfeld, G.~Korchm\'aros and F.~Torres,
\emph{Algebraic Curves over a Finite Field},
Princeton Series in Applied Mathematics,
Princeton University Press, Princeton, 2008.

\bibitem{Roquette}
P.~Roquette,
\emph{Absch\"atzung der Automorphismenanzahl von Funktionenk\"orpern bei Primzahlcharakteristik},
Math. Z. \textbf{117} (1970), 157--163.

\bibitem{Isaacs}
I.~M.~Isaacs,
\emph{Finite Group Theory},
Graduate Studies in Mathematics \textbf{92},
American Mathematical Society, Providence, RI, 2008.

\bibitem{KT}
G.~Korchm\'aros and F.~Torres,
\emph{On the genus of a maximal curve},
Math. Ann. \textbf{323} (2002), 589--608.

\bibitem{Lachaud}
G.~Lachaud,
\emph{Sommes d'Eisenstein et nombre de points de certaines courbes
alg\'ebriques sur les corps finis},
C. R. Acad. Sci. Paris S\'er. I Math. \textbf{305} (1987), no.~16,
729--732.

\bibitem{LiOort}
K.-Z.~Li and F.~Oort,
\emph{Moduli of Supersingular Abelian Varieties},
Lecture Notes in Mathematics \textbf{1680}, Springer, Berlin, 1998.

\bibitem{SGA1}
A.~Grothendieck,
\emph{Rev\^etements \'etales et groupe fondamental (SGA~1)},
Lecture Notes in Mathematics \textbf{224}, Springer, Berlin, 1971.

\bibitem{Stichtenoth}
H.~Stichtenoth,
\emph{Algebraic Function Fields and Codes},
2nd ed., Graduate Texts in Mathematics \textbf{254},
Springer, Berlin, 2009.

\bibitem{Yui}
N.~Yui,
\emph{On the Jacobian varieties of hyperelliptic curves over fields of
characteristic $p>2$},
J. Algebra \textbf{52} (1978), no.~2, 378--410.

\end{thebibliography}

\end{document}
